\section{The arithmetic direction selected by the sweep}
The recent continuations fall into three mathematically different classes.
First, the marked divisor and cofactor correspondences, exact factor-capacity
bounds, and actual-word inverses can produce an integral solution. Second,
SplitZero reconstruction, relation weights, uniform filters, and completed
homotopies protect a computation against loss of information, but their formal
properties alone also hold for prescribed all-failure arrays. Third, the
Leech/monodromy maps have precise retained fibres, but no integral selection
inequality follows just from those inverses. The present continuation pushes
the first class using a coefficient calculation compatible with the second;
it neither discards nor extends the third.

The earlier archive already contains a Star--Kneser target-forcing theorem
\cite[\S12.1, Theorem 12.1, pp.~37--38]{Archive}. Only that local theorem and its
proof are used, not the global conclusion claimed elsewhere in that document.
The new finite comparison also includes short words in \emph{both} residual
and cofactor roles. Comparing only to the preceding sufficient capacity bound
would overstate the computational improvement.

The SplitZero convention is retained:
\[
 G(A)=\{\tau\}\sqcup A^\bullet,\qquad e=0_A^\bullet\ne\tau.
\]
In particular, coefficient reduction can kill an even positive amplitude
without making its original label absent. The scalar $e$ is not an exterior
channel tag, nor a prime multiplicity; those have separately named fields.

\section{Original fixed-seed arithmetic and the two cyclic charts}
Fix
\[
 k\ge3,\quad u=2^{2k-5},\quad m=2^k,\quad o=2^{k-2},\quad
 p\equiv1\pmod8,\quad p>2u,\qquad N=p+4u.                 \tag{A1}
\]
Primality is not needed for the structural theorem. Write the actual integer
factorization as $N=\prod_q q^{e_q}$, with distinct primes $q$.
The factorization is input, not a proposed efficient factorization algorithm.
The original middle states at this fixed seed are in bijection with
\[
 \boxed{Q\mid N,\quad Q\equiv-1\pmod m,\qquad
 R=N/Q,\quad a=(p+R)/4.}                                  \tag{A2}
\]
Indeed $K(u)=2^{k-2}$ and $u\mid a^2$ is equivalent to $K(u)\mid a$.
Since $m\mid4u$, the congruence in (A2) gives $R\equiv-p\pmod m$, hence
$K(u)\mid a$. Both $R$ and $Q$ are $3\pmod4$, so at least three, and
$N<3p$ makes each less than $p$. Also $R\mid p+4u$ implies $R\mid a+u$.
Conversely the original conditions imply (A2). This proves all range and
availability conditions rather than importing mere CRT compatibility.

With $d=\gcd(a,u)$, retain
\[
 (h,r,s_0)=(d^2/u,u/d,a/d),\quad
 \lambda=(r+s_0)/R,\qquad
 (x,y,z)=(a,phs_0\lambda,phr\lambda).                       \tag{A3}
\]
Prime valuations prove integrality of $h,r,s_0$, $a=hrs_0$, $u=hr^2$ and
$\gcd(r,s_0)=1$. As $\gcd(N,u)=1$, all these factors are units modulo $R$;
the gate gives $R\mid r+s_0$. Multiplication of the reciprocal sum by
$phrs_0\lambda$ gives $(p+R)\lambda=4hrs_0\lambda$.
The complete ordered divisor inverse is
\[
 \boxed{u=\frac{pa^2}{Ry-pa}.}                              \tag{A4}
\]
The original exponents in $a$, its centered divisor vector, and the middle
channel travel with (A3). The Type II parametrization is inherited
\cite{ET,TwoBlock}; the forcing calculation below is not attributed to it.

There are two useful cyclic subgroups:
\[
 H_3=\langle3\rangle=\{x:x\equiv1,3\pmod8\},\qquad
 H_{-3}=\langle-3\rangle=\{x:x\equiv1,5\pmod8\},            \tag{A5}
\]
of order $o$ in $(\mathbb Z/m)^\times$. Neither contains $-1$, and both
contain $p$. The identity $v_2(3^{2^j}-1)=j+2$ for $j\ge1$, proved by
successive factorization of a difference of squares, gives the orders;
reduction modulo eight then identifies the two subgroups. Finite logarithms
in either chart are uniquely lifted one binary digit at a time.

Fix $b\in\{3,-3\}$. For every $q\in H_b$ let
$q\equiv b^{a_q}$, $0\le a_q<o$, and let $p\equiv b^s$, $0\le s<o$.
Ignore identity logarithms only in the selected packet. For each remaining
prime define
\[
 w_q=2^{v_2(a_q)},\quad E_q=e_q.
\]
If $s\ne0$, add one \emph{phase item}, with logarithm $-s\bmod o$,
weight $w_p=2^{v_2(s)}$ and capacity one. It is not a prime factor of $N$.
The choice of that item will record residual versus cofactor orientation.
For $q\equiv1\pmod8$, its nonzero weight equals
$2^{v_2(q-1)-2}$; the item is an identity when $v_2(q-1)\ge k$.
The other inside residue class has weight one. These formulas let the
capacity test run without a discrete logarithm.

\section{A complete bounded criterion for reaching the top layer}
\begin{lemma}[Bounded dyadic selection]\label{lem:coins}
Let $o=2^n$, $n\ge1$. Given weights $w_i\in\{1,2,\ldots,2^{n-1}\}$
that are powers of two and capacities $E_i\in\mathbb Z_{\ge0}$, there are
integers $0\le c_i\le E_i$ with
\[
 \sum_i c_iw_i=o-1                                        \tag{B1}
\]
if and only if all the actual prefix capacities satisfy
\[
 \boxed{\sum_{w_i\le2^j}E_iw_i\ge2^{j+1}-1,
 \qquad0\le j<n.}                                        \tag{B2}
\]
The selected items and their original labels can be reconstructed finitely.
\end{lemma}
\begin{proof}
Modulo $2^{j+1}$, the contribution of weights at most $2^j$ must be
$2^{j+1}-1$, so its nonnegative total is at least that amount. This proves
necessity. For sufficiency, aggregate capacities by denomination, keeping a
list of the original labels. Select one unit coin, which (B2) supplies.
Pair all remaining unit coins and replace each pair by a labelled coin of
weight two, leaving at most one unused unit. Divide all remaining weights
by two. The new capacity through denomination $2^j$ equals
\[
 \left\lfloor\frac{C_{j+1}-1}{2}\right\rfloor
 \ge2^{j+1}-1,
 \quad C_{j+1}=\sum_{w_i\le2^{j+1}}E_iw_i.
\]
Induction supplies $2^{n-1}-1$ in the divided system. Expanding each paired
coin gives (B1) with no occurrence used twice. Ties between labels are
resolved in their retained input order.

The executable uses the equivalent descending-denomination greedy selection.
Its exchange proof is as follows. Among coins strictly smaller than a
power of two $w$, any multiset of total at least $w$ contains a submultiset
of total exactly $w$: pair coins upwards; if no weight-$w$ coin is produced,
at most one unpaired coin of each lower weight remains, of total at most
$w-1$. If a representation uses fewer available weight-$w$ coins than the
greedy choice, exchange a lower submultiset of mass $w$ for another such coin.
Iterating preserves a representation and proves the greedy inverse.
\end{proof}

Let $X$ denote a cyclic basis element, not an arithmetic denominator. In
characteristic two the full cyclic coefficient algebra has the explicit
isomorphism
\[
 \mathbb F_2[C_o]=\mathbb F_2[X]/(X^o-1)
 \cong\mathbb F_2[T]/(T^o),\qquad T=X+1,\quad X=T+1.        \tag{B3}
\]
For $a\not\equiv0\pmod o$,
\[
 \operatorname{ord}_T(1+X^a)=2^{v_2(a)}.                   \tag{B4}
\]
To see this, write $a=2^vr$ with $r$ odd and use
$1+(1+T)^a=1+(1+T^{2^v})^r$; its first term is $T^{2^v}$.
All basis coefficients and the inverse Pascal transform are explicit in (B3).
This is an ordinary coefficient-algebra isomorphism; it is not asserted to
be an inverse in the global sparse split semiring.

\begin{theorem}[Actual-factor top-socle packet]\label{thm:top}
Suppose the items of (A5), including the optional phase item, satisfy (B2).
Choose $c_i$ as in (B1). Then
\[
 \boxed{\prod_i(1+X^{a_i})^{c_i}
 =T^{o-1}=1+X+\cdots+X^{o-1}\quad\text{in }\mathbb F_2[C_o].} \tag{B5}
\]
There is an actual unweighted, bounded exponent subfamily whose integer
coefficient at \emph{every} cyclic residue is positive and odd.
\end{theorem}
\begin{proof}
By (B4), the product is $T^{o-1}$ times a unit with constant term one.
All its higher terms vanish modulo $T^o$, proving the first equality.
Repeated squaring gives
$T^{o-1}=\prod_{j=0}^{n-1}(1+X^{2^j})=\sum_{r=0}^{o-1}X^r$.

Write $f\preceq c$ when every binary one-bit of $f$ is also a one-bit of
$c$. Frobenius gives the exact coefficient identity
\[
 (1+Y)^c=\prod_{j:\,c_j=1}(1+Y^{2^j})
 =\sum_{f\preceq c}Y^f\quad\bmod2.                       \tag{B6}
\]
Use the integer packet
\[
 P(X)=(1+X^{-s})^\delta
       \prod_{q:\,c_q>0}\left(\sum_{f_q\preceq c_q}X^{a_qf_q}\right),
 \quad\delta\in\{0,1\}.                                 \tag{B7}
\]
Each original prime exponent appears once and satisfies
$0\le f_q\le c_q\le e_q$. Reduction of (B7) is (B5), so every original
coefficient is odd, hence at least one. Distinct exponent vectors that give
the same residue are still counted separately.
\end{proof}

The group-algebra and binary-cover mechanisms are classical. In particular,
Greene--Higgins \cite[\S1; Theorem 12, p.~7]{GH} describe perfect cyclic
$2$-group covers and their generating polynomials. Here (B7) need not be a
perfect or uniform cover: for logs $(1,3,5)$ and selected capacities
$(3,3,1)$ in $C_8$, its integer coefficients are
\[
 (5,5,3,5,3,3,5,3).
\]
Their parity is one everywhere, but the number of original preimages is not
constant. This is why the positive integer packet, its reduction, and its
full inverse fibres remain separate.

The top line $\mathbb F_2T^{o-1}=\ker(T)$ is nonzero but has augmentation
zero, since $o$ is even. Thus ordinary augmentation is a receiving supported
zero although every original residue is occupied. The result does not infer
positivity from the existence of an abstract socle: it proves that an actual
bounded arithmetic product equals its nonzero top element. Exceeding the top
degree can instead give zero, so the exact equality (B1) is necessary for this
particular proof. No identification with a theta metric is used.

\section{Terminating return to original denominators and counts}
\begin{theorem}[Original witness from the packet]\label{thm:return}
Under (A1), if (B2) holds in either cyclic chart and an actual prime divisor
$q_0\mid N$ lies outside that chart, the fixed-seed original middle branch
contains a positive integral ES witness.
\end{theorem}
\begin{proof}
The element $-q_0^{-1}$ lies in $H_b$; let its logarithm be $r$.
By (B7), choose an actual bounded vector and phase bit with
\[
 Bp^{-\varepsilon}\equiv-q_0^{-1}\pmod m,
 \quad B=\prod_q q^{f_q},\quad f_q\preceq c_q,
 \quad\varepsilon\in\{0,\delta\}.
\]
All prime factors of $B$ are inside the chart and hence different from $q_0$;
therefore $q_0B\mid N$. If $\varepsilon=0$, take $Q=q_0B$ and $R=N/Q$.
If $\varepsilon=1$, take $R=q_0B$ and $Q=N/R$. Their residues are exactly
those of (A2). Equations (A3)--(A4) finish the original return.
\end{proof}

The algorithm multiplies suffixes of the binary factors in (B6) over
$\mathbb F_2[C_o]$. At a target whose current coefficient is odd, exactly one
of the two suffix coefficients for omitting or taking the next factor is odd.
Choose it and update the target. The final target is zero, yielding the
bounded exponent word. An even coefficient on a discarded branch is not
asserted empty. The full fibre of a target is all retained submask vectors
and permitted phase bits satisfying its one modular equation.

For an original returned state with $q_0$, chart and role retained, read
$Q/q_0$ or $R/q_0$ according to the role, then its actual prime valuations.
This gives the complete inverse when it belongs to the selected submasks.
Forgetting the selected chart, outside divisor, or role can identify different
returns. Those are not counted as different original denominator states.
The algorithm is finite in $o$ and the number of binary factors; factoring
$N$ and writing large output integers retain their separate costs.

\paragraph{A genuine lower count.}
Let $M_0=\prod_{q:c_q=0}q^{e_q}$; unallocated occurrences of a partly selected
prime are not included in this \emph{particular disjoint subfamily}.
They still belong to the original full factorization. Define
$A_0=\#\{t\mid M_0:t\notin H_b\}$.
Each such $t$ has a packet word by Theorem~\ref{thm:top}.
Unique factorization makes the return injective within each role, and every
original $Q$ has at most two roles. Therefore
\[
 \boxed{C_k(p)\ge A_0\quad(\delta=0),\qquad
 C_k(p)\ge\lceil A_0/2\rceil\quad(\delta=1).}              \tag{C1}
\]
The exact original mass is
\[
 A_0=\frac12\left(
 \prod_{q^e\parallel M_0}(e+1)-
 \prod_{q^e\parallel M_0}\sum_{j=0}^e\chi_b(q)^j\right),   \tag{C2}
\]
where $\chi_b=1$ inside $H_b$ and $-1$ outside. This expands the actual
divisor list and preserves all multiplicities.

\section{A strict original arithmetic example}
Take the deterministically certified hard prime
\[
 p=852769\equiv169\pmod{840},\quad k=6,\quad u=128,
 \qquad N=853281=3^3\cdot11\cdot13^2\cdot17.              \tag{D1}
\]
In the $3$-chart, $o=16$, the prime logarithms of $3,11,17$ are $1,7,4$,
and $s=8$. The available prefix masses are $(4,4,8,16)$, compared with
$(1,3,7,15)$. Choose $c_3=3$, $c_{11}=0$, $c_{17}=1$, and the phase once.
The actual packet is
\[
 (1+X)(1+X^2)(1+X^4)(1+X^8)=\sum_{r=0}^{15}X^r.           \tag{D2}
\]
Here each target is represented exactly once even before reduction. Its
augmentation modulo two is zero, not absence of those sixteen words.

The outside prime $13$ has target logarithm $11$. The unique packet choice
is $3^3p^{-1}$, so
\[
 R=3^3\cdot13=351,\quad Q=11\cdot13\cdot17=2431,
 \quad(a,h,r,s_0,\lambda)=(213280,8,4,6665,19).
\]
The ordered reciprocal identity is
\[
 \boxed{\frac4{852769}=
 \frac1{213280}+\frac1{863923218520}+\frac1{518483552}.}     \tag{D3}
\]
The original factorization $a=2^5\cdot5\cdot31\cdot43$ gives centered
vector $(2,-1,-1,-1)$ for $u=128$.
The full divisor enumeration has just this one target cofactor.
It has three prime occurrences, while its complementary residual has four;
no prime or two-prime word works in either role.
The preceding $3,11$ paired-gap test requires $e_3\ge6$ here, whereas the
actual exponent is three. Thus (D3) lies beyond both that sufficient bound
and the fair two-occurrence direct baseline.

The methods are not universally ordered. The inherited example
$p=482187176641$, $k=6$, $N=3^7 11^5 37^2$ satisfies the preceding two-block
criterion but fails (B2) in both new charts. Conversely (D1) satisfies the
new criterion but fails that one. Both are retained.
At $p=49681,k=6$, $N=3^3 11\cdot13^2$ the complete branch is empty and
both new charts fail their prefix conditions. These are fixed-seed results,
not assertions about all ES solutions at those primes.

\section{A second retained constraint: concentration at an actual stabilizer}
Let $G=(\mathbb Z/2^k)^\times$ and let $W$ be the set of original divisor
residues of $N$. Assume $-1\notin W$ and let
$H=\{g:gW=W\}$ be its actual stabilizer. Since $1\in W$, $-1\notin H$.
Complementation also excludes $-p$ from $W$. Put
$t_H=1$ if $p\in H$, and $t_H=2$ otherwise.
Then
\[
 \boxed{\sum_{q\bmod m\notin H}e_q\le[G:H]-1-t_H.}         \tag{E1}
\]
For each nontrivial cyclic subgroup $C\le G/H$, there is the additional bound
\[
 \boxed{\sum_{\langle qH\rangle=C}e_q\le |C|-2.}           \tag{E2}
\]
These provide explicit refinements of the archived Star--Kneser interface
\cite{Archive}; Kneser's addition theorem is classical \cite[Theorem 1]{DeVos}.
No global claim of the archived document is an input here.

\begin{proof}
The final divisor set projected to $G/H$ has trivial stabilizer. Adjoin
nonidentity prime occurrences one at a time by $S\mapsto S\cup qS$.
If such a step did not enlarge $S$, it would make $S$ invariant under $q$;
all later products would preserve this invariance, contradicting the final
trivial stabilizer. Starting from one element, the final size is therefore
at least one plus the number of such occurrences. It avoids $t_H$ specified
cosets, proving (E1).

For (E2), isolate occurrences whose projections generate the same cyclic
subgroup $C$. Each is a generator of $C$. Starting with the identity, adjoining
any such generator strictly grows the partial set unless it is already all
of $C$: invariance under a generator would otherwise be impossible.
Thus $|C|-1$ occurrences fill $C$. The full divisor set would then be
$C$-invariant, a contradiction. This proves (E2).
\end{proof}

All possible target-avoiding subgroups have an explicit finite list. With
$n=k-2$, they are the identity and
\[
 \boxed{H_{j,\epsilon}=
 \left\langle(-1)^\epsilon3^{2^{n-j}}\right\rangle,
 \quad1\le j\le n,\quad\epsilon=0,1.}                    \tag{E3}
\]
There are $2k-3$. Projection $G=C_2\times C_{2^n}\to C_{2^n}$ is injective
on a subgroup avoiding $-1$. Such a subgroup is the graph of a homomorphism
from the unique subgroup of order $2^j$ to $C_2$, giving the two choices.
For (E1)--(E2), a candidate on this list is not called the actual stabilizer
unless that has been verified. If every candidate is ruled out, a target is
forced. A surviving candidate alone does not certify emptiness.

A false fixed-seed branch must obey these concentration bounds and, in each
of the two charts with an outside prime, fail at least one (B2). Explicitly,
for some $0\le j<k-2$,
\[
 \sum_{\substack{q\in H_b,\ a_q\ne0\\w_q\le2^j}}e_qw_q
 +\mathbf1_{0<w_p\le2^j}w_p
 \le2^{j+1}-2.                                           \tag{E4}
\]
All other nonidentity inside factors have $q\equiv1\pmod{2^{j+3}}$.
This is a restriction on actual factor multiplicities, not only on the
subgroup they generate. It does not yet contradict all seed factorizations
of one prescribed prime.

\section{Finite comparison, provenance and the remaining arrow}
The complete finite comparison uses the preceding domain $k\ge4$ at all
4519 hard-class primes through two million, with 37289 valid fixed-seed
branches. Direct enumeration finds 4583 original middle states in 3097
occupied branches, occurring at 2411 primes.
\begin{center}
\begin{tabular}{lrr}
\toprule
Test & branches & primes\\\midrule
Preceding $3,11$ capacity &993&990\\
New two-chart top packet &1139&1119\\
Prime/two-prime words, both roles &3051&2387\\
Previous capacity or those short words &3065&2398\\
Same baseline or new packet &3068&2399\\
Full divisor occupancy &3097&2411\\\bottomrule
\end{tabular}
\end{center}
The three strict additional branches occur at
$(p,k)=(852769,6),(1546729,5),(1553281,5)$; only $852769$ is a new prime
for the fair baseline across all its seed branches. The recovered subgroup
criteria alone certify 1753 branches at 1572 primes. These counts describe
sufficient tests, not new overall ES coverage, a least counterexample,
or a universal theorem. The inherited large example above proves that the
observed finite inclusion of the old capacity test in the new one is not
a general implication.

The standalone verifier preserves all original count and inverse data and
checks the selection theorem against bounded knapsack, packet parities
against positive coefficients, and the concentration law against actual
small stabilizers. A second implementation imports none of the first: it
uses direct residue convolution, bounded knapsack, valuation weights, and
explicitly enumerated cyclic subgroups. Reproduction receipts report their
actual scopes; multiple runs are not multiple independent reviews.

The strongest next arrow remains fixed-prime arithmetic. A hypothetical
failure must make every relevant $p+4u$ satisfy at least one specific low-layer
capacity deficiency or have no prime in the required complementary coset,
while also meeting (E1)--(E2). The exact additive differences between these
seed values constrain their factorizations; no theorem here forces their
simultaneous avoidance to be impossible. Nevertheless the top calculation
now proves that an \emph{actually available} bounded product has every target
coefficient nonzero. It does not start from a success projector already
built out of the answer.

No historical-priority claim is made. The group-algebra ingredients, perfect
binary covers, Kneser interface and Type II reconstruction are attributed.
The particular bounded capacity, phase-role return and actual-factor
application are proved here. No Lean build, independent mathematical review,
analytic theta identification or universal ES proof is asserted.
